\section{Background}\label{sec:back}
In ``Patterns for a New Generation: AI and Agents'' \cite{corneli2026patternsnewgenerationai}, we introduced the notion of guiding LLM-based coding agents with a simplified AIF-inspired wrapper.  Those agents still required human guidance to point them at a problem, and the wrapper often gave them a set of hoops that distracted them from the task they had been given.

Motivated by the spatial appearance of the pattern landscape mentioned in Section \ref{sec:observation}, it seemed plausible to build a fuller AIF harness around coding agents that would overcome both failure modes.  The surrounding system would need to select work, check the viability of proposed actions, and learn from outcomes while the coding agents continued to do the coding work.  This is the topic we describe here: the paper is not a sequel but a reboot that treats the definition from the predecessor paper as a working method on several levels.\worknote{IMPROVED: ``select work, check viability, learn from outcomes'' each now carries formal witnesses at pinned records for its constituent equations (belief state, observation, prediction error, forward model, policy set and posterior, precision, policy free energy, depth, temperature, action, Dirichlet accumulation). \wkflag{OUTSTANDING}: the R5 ambiguity estimator (a categorical build is recommended on record), R9 anchor authentication, and the R10 production caller.}

\begin{quote}
\emph{A} design pattern \emph{is a learned generalization arising from a review process that describes how intentions, actions, and outcomes tend to relate across situations.  A candidate pattern becomes a pattern only when it has survived repeated enactment-and-review cycles in which its failure modes are recorded and used to revise it} \cite{corneli2026patternsnewgenerationai}.
\end{quote}

Design patterns play three primary and explicit roles in this work.  First, as Section \ref{sec:observation} suggests, they organise the landscape of possible actions.  Second, they organise the system and its development: the patterns in this paper were written and revised as operational design artifacts while the system was being constructed and tested.  The aim has been to remain as faithful to the active-inference formalism as is realistically possible in this setting.  Third, we express the informal enact--review cycle of the definition above within the active-inference model itself: candidate pattern authoring expands the available policy vocabulary, while witnessed enactment and review decide whether the proposed structure earns standing.

For readers new to active inference: the core idea is that an adaptive agent maintains a \emph{model} of its world, acts to make the world match its preferred expectations, and updates its beliefs from what it then observes --- in a continuous loop.  What distinguishes active inference from a simple optimizer is that action, learning, uncertainty, and preference all live in one loop.  A simple optimization picks the highest immediate payoff; active inference belongs to a more sophisticated class of algorithms that model the value of reducing uncertainty and improving the model itself.  In this paper the framework serves as an operational guide to developing a supplementary harness that adds new \texttt{perceive}, \texttt{believe}, \texttt{evaluate}, \texttt{select}, \texttt{act}, and \texttt{learn} machinery around existing agentic coding systems.  \xglossarynote{}  To earn their status as ``patterns,'' these moves should generalise beyond the current system.  Our first concern, however, has been to get the system running as described and only then investigate those generalisations (some are discussed in \xfuturework{}).

The computational reading of pattern languages explored here predates contemporary LLM agents.  Moran's 1971 review of Alexander, Ishikawa, and Silverstein's multi-service-centre pattern language read the pattern form as a production rule: if a condition is present, then apply an action, because of stated assumptions and evidence \cite{moran1971artificial}.  He also noticed that the cascade of 64 patterns was not just a diagram of contents but a control problem.  Many patterns may be applicable at once; one pattern can enable or invalidate another; and a static wall chart becomes illegible as the language grows.  His proposed Architect's Adviser therefore maintained a model of the evolving project, retrieved applicable patterns, warned about conflicts, asked for missing information, and displayed the because-clause on demand.  This is a direct ancestor of the problem addressed here: pattern cascades need an active, state-aware interpreter whose recommendations remain inspectable.

There is also a bridge between this work and the theoretical side of design pattern research.  Alexander's account of living structure and Salingaros's thermodynamic analogy for architectural ``life'' treat pattern formation as more than a catalogue of reusable forms: they ask how organised structure emerges, stabilises, and becomes coherent across scales \cite{alexander1999a,salingaros1997life}.  Salingaros's measure $L=T\cdot H$\eqanchor{salingaros} combines a temperature-like factor $T$ (detail, curvature, colour) with a harmony factor $H$ (coherence and internal symmetry).  The two vocabularies are related in spirit because both connect order, complexity, evidence, and adaptive structure formation; we do not, however, claim that $L=T\cdot H$ is already a free-energy functional.  The analogy instead provides a route from pattern-language theory to the questions taken up here through active inference: when is adding structure worthwhile, and how can it make future action and perception better organised?
